\documentclass[11pt]{article}
\usepackage[margin=0.9in]{geometry}
\usepackage{amsmath,amssymb,amsthm,booktabs,mathrsfs}
\usepackage[T1]{fontenc}
\usepackage{lmodern,microtype}
\usepackage[colorlinks=true,linkcolor=blue,urlcolor=blue]{hyperref}
\newtheorem{theorem}{Theorem}[section]
\newtheorem{lemma}[theorem]{Lemma}
\newtheorem{proposition}[theorem]{Proposition}
\newtheorem{corollary}[theorem]{Corollary}
\theoremstyle{remark}
\newtheorem{remark}[theorem]{Remark}
\newcommand{\norm}[1]{\left\|#1\right\|}
\title{Toward a four-dimensional Yang--Mills construction\\
\large Notebook 20: the first Wilson small-field integration}
\author{Research working document}
\date{17 October 2026}
\setlength{\emergencystretch}{5em}
\begin{document}
\maketitle
\begin{abstract}
The physical source-coordinate construction of Notebook 19 and
the overlap expansion of Notebook 18 yield an actual first-step
small-field integration of the non-Abelian Wilson action.
For the bare step, every action term is local before and after
the exact pivot substitution. This verifies the support-valued
precision and interaction hypotheses directly. The retained
sharp cutoff is fixed; its nonlinear Jacobian, Haar density,
gauge remainder and normalization terms are kept explicitly.
The normalized fluctuation logarithm has connected activities
whose anchored norm is bounded by
$C[gR^3+\exp(-cR^2)]$, for a logarithmic cutoff $R$.
Constants are uniform in periodic volume at fixed blocking
ratio. Strict support here is in the independent current
background. Its transport to the next coarse field, the
subsequent quantum iterations, and the complementary
large-field regions remain the next construction steps.
\end{abstract}

This notebook applies the internally checked finite-regulator
results in Notebooks 15--19. It is a proof of a specified
small-field contribution. Completing the full existence
argument requires that contribution to be incorporated into
an iteration that controls the whole regulator measure and
the continuum observable hierarchy.


% No preamble. Actual bare first-step application. Unique prefix bs-.
\section{An actual first-step Yang--Mills small-field integral}
\label{bs-section}

We now set the preceding depth to $M=1$. The current classical
action $S_1$ is the original, coupling-free Wilson action. This
specialization supplies the interaction and precision-locality
hypotheses directly: no localization theorem for a previously
minimized action is needed. The block side $L$ is fixed throughout.
The conclusion is a convergent representation of the specified
first-step small-field integral, with constants uniform in periodic
volume. It does not include the complement of that integral or
an iteration of quantum effective interactions.

Use the exact ordinary constraint, pivot chart $\Phi_V$, and
absolute contour gauge function $\mathcal G$ of Notebook 19.
Set $a=\log V/i$ and write
\begin{equation}\label{bs-action}
 \mathcal F_V(x)=S_1\bigl(\log(e^{i\Phi_V(x)}V)/i\bigr)
                   +\mathcal G(e^{i\Phi_V(x)}V).
\end{equation}
This symbol is the scalar action plus gauge function; it is not
the group-valued coarse constraint. On the physical family
$V=V_*(W)$ constructed in Notebook 19,
\begin{equation}\label{bs-centered}
 \mathcal F_V(0)=S_1(a),\qquad D\mathcal F_V(0)=0,\qquad
 D^2\mathcal F_V(0)=A(V).
\end{equation}
On a full product domain of independent small complex current
link potentials, define $A(V)=D^2\mathcal F_V(0)$ using the
actual absolute gauge function. This extension need not be
stationary. Its flat value is
$A_0=K_L^T(q_1+T_L^TT_L)K_L$ and has the spectral and
exponential inverse bounds of Notebooks 17 and 19.

\subsection{Finite local action and precision labels}

Assign retained positive bonds to their $L$-blocks. There are
$d(L^d-1)$ retained bonds per block in the stated periodic pivot
geometry, so the number of real fluctuation coordinates per
cell is the fixed integer $D=d(L^d-1)\dim\mathfrak g$.
The background coordinates assigned to a cell include its pivot
bonds as well. These background and fluctuation coordinate
counts need not be equal. All supports below refer to these cells.

Each Wilson plaquette reads four current links. Substitution of
$e^{i\Phi_V(x)}V$ enlarges this to a fixed neighborhood because
each pivot correction reads one fixed local stencil. Each gauge
contour likewise reads a fixed neighborhood. Hence there is an
exact sum
\begin{equation}\label{bs-local-action}
 \mathcal F_V(x)=\sum_r f_r(V|_{Y_r},x|_{Y_r}),
\end{equation}
with at most $s_L$ cells in each $Y_r$, tree length at most
$\ell_L$, and at most $n_L$ labels containing any given cell.
Periodic occurrences are counted with their multiplicities.
All $f_r$ are holomorphic and uniformly bounded on a fixed
larger background/fluctuation polydisc. The constants depend
on $L$ and the fixed gauge group, not on volume.

Local Cauchy estimates therefore prove, on a smaller polydisc,
\begin{align}
 A(V)&=A_0+\sum_r B_r(V|_{Y_r}),\label{bs-precision}\\
 B_r(V)&=D_x^2f_r(V,0)-D_x^2f_r(1,0),\notag\\
 \nu_t(b)&:=\sup_i\sum_{r:i\in Y_r}\sum_{j\in Y_r}
 e^{t\ell(Y_r)}\|(B_r)_{ij}\|_{\|a\|_\infty<b}
 \le C_{L,t}b.\label{bs-precision-norm}
\end{align}
The matrices $B_r$ are symmetric and supported in rows and
columns in $Y_r$; they read only the background there.
This is a support-valued precision decomposition, not merely
a bound on the summed kernel. It retains every constraint and
ordinary-coordinate curvature term automatically because it
differentiates the full local functions in \eqref{bs-local-action}.

Define the third-and-higher action interaction on the product
domain by
\begin{equation}\label{bs-action-remainder}
 v_r^{\rm act}(V;g,z)=-g^{-2}\left[
 f_r(V,gz)-f_r(V,0)-gD_xf_r(V,0)z
       -\tfrac12g^2D_x^2f_r(V,0)[z,z]\right].
\end{equation}
Its singularity at $g=0$ is removable. For $R\ge1$ and
$|g|R$ below a fixed radius, local Taylor estimates give
\begin{equation}\label{bs-remainder-norm}
 \sup_i\sum_{r:i\in Y_r} e^{a_0|Y_r|+\kappa\ell(Y_r)}
       \|v_r^{\rm act}\|_{b,R}
       \le C_{L,a_0,\kappa}|g|R^3.
\end{equation}
Here the norm includes the independent background product ball
and the complex fluctuation polydisc of radius $R$ in each
retained bond coordinate. One obtains the bound by applying
one-variable Taylor's formula to $t\mapsto f_r(V,tz)$;
the number of variables per label is bounded at fixed $L$.
There is no sum over arbitrarily many fine variables: this
section is specifically the bare $M=1$ step.

Add the normalized ambient straightening-Jacobian and Haar
logarithms from Notebook 19, grouping all primitive terms with
equal support. Denote the resulting interaction by $v_Y$.
The gauge remainder is already part of
\eqref{bs-action-remainder} and is not added twice. For any
fixed $a_0,\kappa$,
\begin{equation}\label{bs-full-interaction}
 \sup_i\sum_{Y\ni i} e^{a_0|Y|+\kappa\ell(Y)}
           \|v_Y\|_{b,R}
 \le C_{L,a_0,\kappa}
       (|g|R^3+|g|R+|g|^2R^2)
 \le C'_{L,a_0,\kappa}|g|R^3,
\end{equation}
where $|g|R\le1$. The same bound, up to a fixed factor,
holds for $\|e^{v_Y}-1\|$ when the displayed quantity is
small. Every term is strictly local in the independent current
background, even though substituting $V_*(W)$ need not preserve
finite support in $W$.

\subsection{The retained source cutoff fits the cell theorem}

Let
\begin{equation}\label{bs-cutoff}
 \chi_R(z)=\prod_{b\ {\rm retained}}\mathbf1_{\{|z(b)|<R\}},
\end{equation}
where $|\cdot|$ is the real color norm. In each cell let $\chi_i$
be its product of these bond indicators. Then
$1-\chi_i\le\mathbf1_{\{\|z_i\|_2\ge R\}}$.
Consequently the cell tail estimates of Notebook 18 also apply
to these rectangles: their proofs only use a fixed real
indicator, bounded by one, and the stated bound on its bad
event. The interaction supremum is taken on the actual
retained-bond rectangle. Fixed factors converting bond and
cell norms are included in the constants in
\eqref{bs-full-interaction}. No complex indicator is introduced.
The cutoff is an exact pullback because
$\Phi_V(gz)(b)=gz(b)$ on every retained bond.

\subsection{Closing all smallness conditions in the first step}

Put $C_0=A_0^{-1}$. Choose $0<\kappa<m$ with
$\kappa+2m$ below its established inverse-decay exponent and
set $K=\|C_0\|_{\kappa+2m}$. Its real spectral bounds are
$\lambda=B_*^{-1}$ and $\Lambda_0=\gamma_*^{-1}$.
Fix $0<\theta\le1/4$, for example $\theta=1/(4D)$.
The constants in Notebook 18's tunable strictly local theorem
are now fixed independently of volume and of the coupling.
Equation \eqref{bs-precision-norm} permits a positive choice
of $b$ satisfying all its precision smallness conditions.
Indeed its $\nu$ tends to zero with $b$, and all other
constants at this stage are finite and fixed.

Let $R(g)=c(\log g^{-2})^p$, where $c>0$, $p>1/2$,
and $g>0$ is small. Take $a_0=4+\log\mathcal K_\theta$
in \eqref{bs-full-interaction}, with the density constant of
that theorem. Its interaction and bad-cell parameters obey
\begin{equation}\label{bs-parameters}
 \epsilon(g)\le C_LgR(g)^3,\qquad
 \tau(g)\le C_L e^{-c_LR(g)^2}.
\end{equation}
Both tend to zero. The recovered whole field stays inside
the chart since $\|\Phi_V(gz)\|_\infty\le C_LgR(g)$.
Thus a further, volume-independent coupling ceiling enforces
the chart condition and both quarter budgets of Notebook 18.
The choice of $b$ precedes that ceiling; small coupling has
not been used to repair a failed precision estimate.

\begin{theorem}[The bare retained small-field integral]
\label{bs-theorem}
There are positive $b,g_0,a_*,C,c_*$, independent of compatible
periodic volume, such that for $0<g<g_0$ the function
\begin{equation}\label{bs-normalized-integral}
 Z_g(V)=\int d\mu_{A(V)^{-1}}(z)\,
                  \chi_{R(g)}(z)\exp\!\left(\sum_Y v_Y(V,z)\right)
\end{equation}
has a holomorphic logarithm on the independent product ball
$\|\log V/i\|_\infty<b$, and a convergent expansion
\begin{equation}\label{bs-local-log}
 \log Z_g(V)=\sum_H E_{g,H}(V|_H),\qquad
 \sup_i\sum_{H\ni i}e^{a_*|H|}\|E_{g,H}\|_b
 \le C\{gR(g)^3+e^{-c_*R(g)^2}\}.
\end{equation}
The supports $H$ are connected periodic cell sets. The same
canonical expansion differs from that for the pure Gaussian
cutoff by at most $CgR(g)^3$ in this norm.
\end{theorem}
\begin{proof}
The covariance baseline, support-valued precision, product
domain, bounded local interaction multipliers and retained
cutoff have all been checked above. Apply the tunable strictly
local theorem and its periodic extension from Notebook 18.
The right side follows from \eqref{bs-parameters}. The
analytic continuation is that of the fixed real Gaussian
integral, with its determinant branch chosen at the real
positive baseline. It is not a distributional complex
constraint integral.
\end{proof}

The product-domain extension subtracts the local linear
terms in \eqref{bs-action-remainder}. On the physical
stationary family their sum is zero by \eqref{bs-centered},
so \eqref{bs-normalized-integral} is exactly the normalized
fluctuation factor of the original action plus gauge.
At nonstationary backgrounds it is the explicitly defined
centered extension, not the unmodified Wilson integral.

\subsection{The normalization terms also have local expansions}

The delta-elimination factor is
\begin{equation}\label{bs-delta-factor}
 \mathcal D(V)=\prod_c\det(A_{V,p(c)})^{-1},\qquad
 \log\frac{\mathcal D(V)}{\mathcal D(1)}
         =-\sum_c\operatorname{Tr}\log\pi_c(V).
\end{equation}
It is a sum of fixed-range local $O(b)$ terms. It has no
coupling suppression and is not absorbed into
\eqref{bs-local-log}. The Gaussian normalization similarly
contributes
\begin{equation}\label{bs-gaussian-det}
 \mathcal Q(V)=-\tfrac12\log\det(A(V)A_0^{-1})
 =\tfrac12\sum_{n\ge1}\frac{(-1)^n}{n}
                 \operatorname{Tr}[(C_0 B(V))^n],
 \qquad B(V)=\sum_r B_r(V).
\end{equation}
The branch is zero at $V=1$.

For completeness, this determinant has a strict local
expansion with volume-independent anchored $O(b)$ norm.
Fix a small $t>0$ below the covariance exponent. A term
in the expanded trace is a closed chain of covariance
links and precision labels. Realize each covariance link
by a chosen shortest periodic path and each label by a
chosen minimal tree on $Y_r$. Their union $H$ is connected
and contains every background coordinate read by the term.
If $i$ is the initial trace cell and $\ell$ is the sum of
these path/tree lengths, then
\begin{equation}\label{bs-hull-estimate}
 |H|\le1+\ell,\qquad d(i,j)\le\ell\quad(j\in H).
\end{equation}
Writing $K_t=\|C_0\|_t$ and using
\eqref{bs-precision-norm}, the sum of absolute matrix
products at order $n$, rooted at $i$ and weighted by
$e^{t\ell}$, is at most $D(K_t\nu_t(b))^n$.
This follows successively from each weighted covariance
row sum and then each weighted precision-label row sum;
closing the chain only decreases the unrestricted endpoint
sum. Choose $a_d,\beta>0$ with $a_d+\beta\le t$.
For every $j\in H$, \eqref{bs-hull-estimate} gives
$e^{a_d|H|+\beta d(i,j)}\le e^{a_d}e^{t\ell}$.
Summing first over rooted chains and then their root yields
\begin{equation}\label{bs-det-bound}
 \sup_j\sum_{H\ni j}e^{a_d|H|}\|\mathcal Q_H\|_b
 \le\frac{De^{a_d}}2
       \left(\sum_i e^{-\beta d(i,j)}\right)
       [-\log(1-K_t\nu_t(b))]\le C_L b.
\end{equation}
The periodic lattice sum is bounded by its infinite-lattice
counterpart. A further decrease of $b$ ensures
$K_t\nu_t(b)<1$. Absolute convergence justifies grouping by
$H$ and holomorphy of each local activity. The local delta
terms in \eqref{bs-delta-factor} meet the same type of bound.

\subsection{Exact relation to the constrained group integral}

This subsection makes the field-independent constants and the
small-field restriction explicit. Let $q=\dim\mathfrak g$,
let $N_b$ be the number of current positive bonds, $N_c$ the
number of next-coarse positive bonds, and set
$n=q(N_b-N_c)$. Choose Lebesgue coordinates orthonormal for
the normalized trace and write normalized Haar measure near
the identity as
$d_H(e^{iX})=c_H j(X)\,dX$, where $j(0)=1$.
A normalized group delta at the identity consequently pulls
back to $c_H^{-1}\delta(X)$.

For real physical $V=V_*(W)$ define $I_g^{\rm sf}(V)$ as
the group integral of
$\exp[-g^{-2}(S_1(U)+\mathcal G(U))]$
with constraint $\mathcal M_L(U)=W$, restricted to the local
chart $U=e^{i\Phi_V(gz)}V$ and the retained rectangle
\eqref{bs-cutoff}. There are no other indicators in this
definition. The exact ambient change of variables and pivot
elimination from Notebook 19 give
\begin{equation}\label{bs-exact-group-integral}
 \begin{split}
 I_g^{\rm sf}(V)={}&
 c_H^{N_b-N_c}g^n\mathcal D(1)
       (2\pi)^{n/2}(\det A_0)^{-1/2}\,
       e^{-S_1(a)/g^2}\\
 &\times\exp\left\{
 \log\frac{\mathcal D(V)}{\mathcal D(1)}
               +\mathcal Q(V)+\sum_HE_{g,H}(V|_H)\right\}.
 \end{split}
\end{equation}
All terms on the second line have convergent strictly local
current-background expansions. The classical term on the
first line is kept explicit. At the physical stationary
background it equals the next constrained minimum action by
the minimum comparison in Notebook 19.

The exponential gauge can be inserted into a gauge-invariant
constrained integral on a rooted-gauge-invariant admissible
domain where the contour mean and ordinary average are defined
and have their stated equivariance. This domain condition is
needed for the source's logarithmic averaging maps; no global
definition on every group configuration is assumed. The
insertion has a known scalar normalization. Let
\begin{equation}\label{bs-gauge-normalization}
 z_G(g)=\int_G d_Hu\,
       \exp[-g^{-2}(1-\operatorname{Re}\tau u)],
 \qquad N_s=\#\{\hbox{noncenter current sites}\}.
\end{equation}
For rooted site gauges, contour equivariance gives
$P_{U^u}(y,x)=P_U(y,x)u(x)^{-1}$.
Haar invariance thus makes the gauge average of
$e^{-g^{-2}\mathcal G(U^u)}$ exactly $z_G(g)^{N_s}$,
independently of $U$ in that domain. Averaging a
gauge-invariant constrained integral over that domain proves
the gauge-insertion identity with factor
$z_G(g)^{-N_s}$. This identity holds before the chart
restriction; the chart piece in
\eqref{bs-exact-group-integral} then has that same external
factor. It does not equate a nongauge-invariant sharp chart
piece before and after gauge insertion.

\subsection{What this first step establishes}

This is an actual non-Abelian Wilson small-field integration,
including the source constraint geometry, fluctuation
precision, retained cutoff, classical remainder, Haar density,
straightening Jacobian, delta determinant and Gaussian
determinant. Its bounds are uniform in periodic volume.
It closes the first bare small-field step with no unknown
coefficient hypotheses.

Three further passages are distinct. First, composing the
strictly local current-background activities with $V_*(W)$
requires a locality theorem in the new coarse background.
Second, subsequent steps must control the previously generated
quantum interactions with uniform bounds. Notebook 21 supplies
the intervening localization of the minimized classical action
and one integration of it uniformly through earlier classical
depth. Third, the regulator construction must
control the complement of the chart, including compatible
large-field histories, and produce the observable bounds
needed for continuum reconstruction. These are not inferred
from the first-step theorem.

\begin{thebibliography}{9}
\bibitem{prior20}
\emph{Toward a four-dimensional Yang--Mills construction},
research Notebooks 15--19 in this workspace.
Notebook 18 proves the overlapping-interaction integration
theorem and its strict support and periodic extensions.
Notebook 19 gives the exact physical pivot chart, background
Gaussian and local measure factors used here.
\bibitem{balaban10920}
T.~Ba\l aban,
\emph{Renormalization group approach to lattice gauge field
theories I}, Commun. Math. Phys. \textbf{109} (1987), 249--301,
\href{https://scispace.com/pdf/renormalization-group-approach-to-lattice-gauge-field-3ygai0vxt1.pdf}{original article copy}.
\bibitem{balaban11620}
T.~Ba\l aban,
\emph{Renormalization group approach to lattice gauge field
theories. II. Cluster expansions},
Commun. Math. Phys. \textbf{116} (1988), 1--22,
\href{https://scispace.com/pdf/renormalization-group-approach-to-lattice-gauge-field-4penm4k7hc.pdf}{original article copy}.
\bibitem{source20}
E.~Sussman and collaborators,
\href{https://github.com/ethanwsussman/Yang--Mills-LLM-experiment}{Yang--Mills-LLM-experiment},
deposited raw accounts at commit
\texttt{e392ef4501\allowbreak 3a581664ba\allowbreak b76b908671\allowbreak 70b30fe10e}.
\end{thebibliography}
\end{document}
